
\section{A computer-assisted proof of renormalization hyperbolicity}
\label{sec:pfmain1}

\subsection{Contraction Mapping Principle}

Our proof of the existence of a fixed point $\zeta^*=(\eta_*,\xi_*)$ is a computer implementation of the Contraction Mapping Principle for the Newton map for the operator $\cRG$ in the space  $\cA_1(U,V)$.


Let $\cF$ be an operator analytic in some neighborhood $\cN$ in a Banach space $\cZ$, which has a fixed point $\zeta_*$ in this space,
 and whose differential $D|_{\zeta_*} \cF$ is a hyperbolic linear map in $\cN$.  Suppose, that one knows an approximation
of the fixed point  $\zeta_0 \approx \zeta_*$.
 Set
$$M \equiv \left[ \field{I}-L_0 \right]^{-1},$$
where $L_0$ is an approximation of $D \cF[\zeta_0]$. For all $\zeta$, such that $\zeta_0+M \zeta \in \cN$, set
$$N_{\cF,L_0}[\zeta]=\zeta+\cF[\zeta_0+M \zeta]-(\zeta_0+M \zeta).$$


Notice, that if $\hat\zeta$ is a fixed point of $N_{\cF,L_0}$, then $\zeta^*=\zeta_0+M \hat\zeta$ is a fixed point of $\cF$.    


The operator $\field{I}-L_0$ is indeed invertible if $L_0$ is close to $D \cF[\zeta_0]$, since $D \cF$ is hyperbolic at $\zeta_0$. If $\zeta_0$ is a reasonably good approximation of the true fixed point of $\cF$, then the operator $N_{\cF,L_0}$ is expected to be a strong contraction in a neighborhood of $0$:
\begin{eqnarray}
\nonumber D N_{\cF,L_0}[\zeta] &=& \field{I} +D \cF[\zeta_0+M \zeta] \cdot M -M= \left[ M^{-1} +D \cF[\zeta_0+M \zeta]  -\field{I} \right] \cdot M \\
\nonumber &=& \left[\field{I}-L_0  +D \cF [\zeta_0+M \zeta]  -\field{I} \right] \cdot M = \left[ D \cF [\zeta_0+M \zeta] - L_0 \right] \cdot M.
\end{eqnarray}

The last expression is typically small in a small neighborhood of $0$, if the norm of $M$ is not too large (if $M$ is large, one might have to find a better approximation $\zeta_0$ and take a smaller neighborhood of $0$).  The following well-known version of the contraction mapping theorem specifies a sufficient condition for the existence of a fixed point $\hat\zeta$:


\begin{cmp}
Let $\cN$ be an open subset of a Banach space $\cZ$.
Suppose that  $N_{\cF,L_0}$ is a well-defined  analytic operator from  $\cN$ to $\cZ$. Let $\zeta_0\in \cN$ and $B_\delta(\zeta_0) \subset \cN$ be such that
$$\| D N_{\cF,L_0}[\zeta]\| \le  \cD <1, $$
for any $\zeta \in  B_\delta(\zeta_0)$, and
$$\|N_{\cF,L_0}[\zeta_0]-\zeta_0\|\le \epsilon.$$

If $\epsilon < (1-\cD) \delta$ then the operator $N_{\cF,L_0}$ has a fixed point $\hat\zeta$ in $B_\delta(\zeta_0)$, such that
$$\| \hat\zeta-\zeta_0\| \le {\epsilon \over 1-\cD }.$$
\end{cmp}

In practice, for a computer-assisted proof, the linear map $L_0$ is always chosen as a finite-dimensional approximation of $D \cF[\zeta_0]$.

\subsection{Differential of the renormalization operator}

We will now compute the action of the differential $D \cRG(\phi,\psi)$ on a vector $(u,v) \in T_{(\phi,\psi)} \cA_1(U,V)$. First, 
let $a$, $b$, and $c$ be the coefficients of the projection $\cP$ considered as functions of the pair $(\phi,\psi)$. Then the differential
%%denote $D c (u,v)$, $D a (u,v)$ and  $D b (u,v)$ actions of derivatives of $c$, $a$, $b$ as functions of $(\phi,\psi)$ on a vector $(u,v)$, then 
\begin{equation}
D \cP(\phi,\psi) (u,v)=(u,v)+( Da (u,v) z^2+Db (u,v) z^3, Dc (u,v)). 
\end{equation}
The differential $D c (u,v)$ can be found easily from the equation $(\ref{eq:newFs1})$:
$$D c(u,v)=-v(0),$$
while   $D a (u,v)$ and  $D b (u,v)$ are found by differentiating the system $(\ref{eq:newFs2})$:
\begin{eqnarray} \label{eq:newFs3}
A \cdot \left[
\begin{array}{c}
a \\
b
\end{array}
\right]+B \cdot \left[
\begin{array}{c}
D a (u,v) \\
D b (u,v)
\end{array}
\right]=\left[
\begin{array}{c}
c_1 \\
c_2
\end{array}
\right],
\end{eqnarray}
where $$A=\left[ a_{11} \  a_{12}  \atop a_{21}  \  a_{22} \right]$$ with
\begin{eqnarray}
\nonumber a_{11}&=&4 \hat{\psi}(0)^{3}(v(0)+ Dc (u,v))=0, \\
\nonumber a_{12}&=&6 \hat{\psi}(0)^{5}(v(0)+ Dc (u,v))=0, \\
\nonumber a_{21}&=&6 \hat{\psi}(0)^{2}  \hat{\psi}'(0) (v(0)+ Dc (u,v))+ 2 \hat{\psi}(0)^{3} v'(0)=2 \hat{\psi}(0)^{3} v'(0), \\ 
\nonumber a_{22}&=&15 \hat{\psi}(0)^{4} \hat{\psi}'(0)(v(0)+ Dc (u,v))+3 \hat{\psi}(0)^{5} v'(0)= 3 \hat{\psi}(0)^{5} v'(0) \Rightarrow \\
\nonumber A &=&\left[
\begin{array}{c c}
0 \quad  & \quad 0   \\
2 \hat{\psi}(0)^{3} v'(0) \quad  & \quad 3 \hat{\psi}(0)^{5} v'(0)
\end{array}
\right], \\
\nonumber  B&=&  \left[
\begin{array}{c c}
\hat{\psi}(0)^{4} \quad  & \quad \hat{\psi}(0)^{6}   \\
 2 \hat{\psi}(0)^{3}  \hat{\psi}'(0) \quad  & \quad  3 \hat{\psi}(0)^{5} \hat{\psi}'(0)
\end{array}
\right],\\
\nonumber c_1&=&(v(\phi(0)^2)+ Dc (u,v))+2 \hat{\psi}'(\phi(0)^2) \phi(0) u(0) \\
\nonumber &\phantom{=}& \phantom{(v(\phi(0)^2)}- u(\hat{\psi}(0)^2)-2 \phi'(\hat{\psi}(0)^2)\hat{\psi}(0) (v(0)+ Dc (u,v)) \\
 \nonumber &=&(v(\phi(0)^2)+ Dc (u,v))+2 \hat{\psi}'(\phi(0)^2) \phi(0) u(0) - u(\hat{\psi}(0)^2)
\end{eqnarray}
\begin{eqnarray} \label{eq:matrices}
\nonumber c_2&=&v'(\phi(0)^2) \phi(0) \phi'(0)+2 \hat{\psi}''(\phi(0)^2) \phi(0)^2 \phi'(0) u(0)+\hat{\psi}'(\phi(0)^2) u(0) \phi'(0)+\\
\nonumber &\phantom{=}& \phantom{v'(\phi(0)^2) \phi(0) \phi'(0)}+ \hat{\psi}'(\phi(0)^2) \phi(0) u'(0)-u'(\hat{\psi}(0)^2) \hat{\psi}(0) \hat{\psi}'(0) -\\
\nonumber &\phantom{=}& \phantom{v'(\phi(0)^2) \phi(0) \phi'(0)}-2 \phi''(\hat{\psi}(0)^2) \hat{\psi}(0)^2 \hat{\psi}'(0) (v(0)+ Dc (u,v)) \\
\nonumber &\phantom{=}&\phantom{v'(\phi(0)^2) \phi(0) \phi'(0)}-\phi'(\hat{\psi}(0)^2) (v(0)+ Dc (u,v)) \hat{\psi}'(0)-\phi'(\hat{\psi}(0)^2) \hat{\psi}(0) v'(0) \\
\nonumber &=&v'(\phi(0)^2) \phi(0) \phi'(0)+2 \hat{\psi}''(\phi(0)^2) \phi(0)^2 \phi'(0) u(0)+\hat{\psi}'(\phi(0)^2) u(0) \phi'(0)+ \\
\nonumber &\phantom{=}& \phantom{v'(\phi(0)^2) \phi(0) \phi'(0)} +  \hat{\psi}'(\phi(0)^2) \phi(0) u'(0)-u'(\hat{\psi}(0)^2) \hat{\psi}(0) \hat{\psi}'(0)- \\
\nonumber &\phantom{=}& \phantom{v'(\phi(0)^2) \phi(0) \phi'(0)} - \phi'(\hat{\psi}(0)^2) \hat{\psi}(0) v'(0) \\
\nonumber &=&\left(v'(\phi(0)^2) +2 \hat{\psi}''(\phi(0)^2) \phi(0) u(0) \right)\phi(0) \phi'(0)  + \\
\nonumber  &\phantom{=}&\phantom{\left(v'(\phi(0)^2) \right.}   + \hat{\psi}'(\phi(0)^2) \left(u(0) \phi'(0)+\phi(0) u'(0) \right)-\\
\nonumber  &\phantom{=}&\phantom{\left(v'(\phi(0)^2) \right.}  - u'(\hat{\psi}(0)^2) \hat{\psi}(0) \hat{\psi}'(0)-\phi'(\hat{\psi}(0)^2) \hat{\psi}(0) v'(0).
\end{eqnarray}
The action of the differential of the operator $\cRG$ on a vector $(u,v)$, therefore, is given by
$$D \cRG (\phi,\psi) (u,v)=(\tilde{u},\tilde{v})+( Da (\tilde{u},\tilde{v}) z^2+Db (\tilde{u},\tilde{v}) z^3, -\tilde{v}(0)),$$
where  $Da (\tilde{u},\tilde{v})$ and  $Db (\tilde{u},\tilde{v})$ are the solutions of the system $(\ref{eq:newFs3})$ for the functions  $\tilde{u}$ and $\tilde{v}$ given by
\begin{eqnarray}
\nonumber \tilde{u}&:=&\Pi_\phi D \cR(\phi,\psi) (u,v)\\
\nonumber &=& c \circ \left\{\left( \left(\lambda^{-1} \circ \phi \circ q_2 \right)' \circ \psi \circ \lambda^2 \right) \cdot v \circ  \lambda^2 + \lambda^{-1} \circ u \circ q_2 \circ \psi \circ \lambda^2 -\right.\\
\nonumber  &\phantom{=}& \phantom{c \circ \left.\left( \left(\lambda^{-1} \circ \phi \circ q_2 \right)' \circ \psi \circ \lambda^2 \right) \cdot v \circ  \lambda^2  \ \right.} -  {v(0) \over \lambda^2} \circ \phi \circ q_2 \circ \psi \circ \lambda^2  + \\
\nonumber  &\phantom{=}& \phantom{c \circ \left.\left( \left(\lambda^{-1} \circ \phi \circ q_2 \right)' \circ \psi \circ \lambda^2 \right) \cdot v \circ  \lambda^2 \  \right. }  \left. + 2 \left(\left( \phi \circ q_2 \circ \psi \right)' \circ \lambda^2  \right) \cdot D \la v \right\} \circ c, \\
\nonumber \tilde{v}&:=&\Pi_\psi   D \cR(\phi,\psi) (u,v)= c \circ \left\{ \lambda^{-1} \circ u \circ \lambda^2+{v(0) \over \lambda^2} \circ \phi \circ \lambda^2 +\left( \phi' \circ \lambda^2\right) \cdot D \la v \right\}  \circ c,
\end{eqnarray}
here $\Pi_{\phi/\psi}$ denote the projections on the corresponding components of the pair, and  $D \la v(z):=v(0) z$.


We remark, that the action $D \cR(\phi,\psi)(u,v)$ is of the form $c \circ L (u,v) \circ c$, where $L$ is a linear map and $c$ is the complex conjugation. Therefore, the differential $ D \cR(\phi,\psi)$ is an {{\it anti-linear}} map.



\subsection{Existence of the renormalization fixed point}

A direct application of the Contraction Mapping Principle, in the form stated above, with the operator $\cRG$ acting on $\cA_1(U,V)$ is impossible: the map $\I - L_0$, being neither linear nor anti-linear on $T_{(\phi,\psi)}\cA_1(U,V)$, does not admit a matrix representation in $GL(N,\C)$ (here, $L_0$ is assumed to be an approximation of the action of $D \cRG(\phi,\psi)$ on some $N$-dimensional subspace of  $T_{(\phi,\psi)}\cA_1(U,V)$).

One way to circumvent this problem is to consider the linear operator $D \cRG \circ c$, where $c$ is complex conjugation. We, however, take a different approach, and think of $\cA_1(U,V)$ as a {\it real} Banach space, and expand its elements into
%%in the Banach  space $\tilde{\cH}(U,V) \approx \cH(U,V)$ 
{\it real-valued} vectors over the basis
\begin{eqnarray}
\nonumber \cL:&=&\left\{ \{ (l_\phi^k,0) \}_{k=0}^\infty,\{(0,l_\psi^k)\}_{k=0}^\infty,\{(i l_\phi^k,0)\}_{k=0}^\infty,\{(0,i l_\psi^k)\}_{k=0}^\infty \right\}\\
\nonumber &=&:\left\{ \{ u_k \}_{k=0}^\infty,\{v_k\}_{k=0}^\infty,\{\upsilon_k\}_{k=0}^\infty,\{w_k\}_{k=0}^\infty \right\},
\end{eqnarray}
where 
$$l_\phi(z)={z-c_\phi \over r_\phi}, \quad l_\psi(z)={z-c_\psi \over r_\psi}.$$
%%equipped with the $l_1$-norm $(\ref{eq:l1norm})$ with respect to which the space is complete.

This real Banach space will be referred to as  $\tilde{\cA}_1(U,V)$.

Of course, the space  $\tilde{\cA}_1(U,V)$ has a linear complex structure defined on it, specifically, 
$$J(u_k)= \upsilon_k, \  J(\upsilon_k)= -u_k, \ J(v_k)= w_k, \  J(w_k)= -\upsilon_k,$$
which, however, will play no role in our computation of the differential.

Notice, that by definition, balls $B_r((\phi,\psi))$ in both spaces   $\tilde{\cA}_1(U,V)$ and  $\cA_1(U,V)$ coincide.

In this setting, an approximation $L$  of the action of the  differential $D \cRG(\phi,\psi)$ on an $N$-dimensional subspace of  $T_{(\phi,\psi)} \tilde{\cA}_1(U,V)$ admits a representation in $GL(N,\R)$.

To construct an approximation of the differential of renormalization, we first find an  approximate renormalization fixed point by renormalizing several times the pair $(\id,p_\theta)$, where $p_\theta$ is the conjugate $ l_\psi \circ P_\theta \circ l_\psi^{-1} $ of the quadratic polynomial $P_\theta=e^{2 \pi i \theta} z  -0.5 e^{2 \pi i \theta} z^2$, $\theta=(\sqrt{5}-1)/2$, by the linear transformation $l_\psi$.  The resultant pair is referred to as $(\phi_a,\psi_a)$.

A linear approximation $L_a$ of the action of the renormalization differential  $D \cRG(\phi_a,\psi_a)$ is constructed on the span
$${\rm span} \left\{ \{ (l_\phi^k,0) \}_{k=0}^{N_1},\{(0,l_\psi^k)\}_{k=0}^{N_2},\{(i l_\phi^k,0)\}_{k=0}^{N_1},\{(0,i l_\psi^k)\}_{k=0}^{N_2} \right\},$$
where 
$$N_1=120, \quad N_2=160,$$
and is of the form 
$$L_a=\left[ \begin{array}{c c c c}
M_{1 1} & M_{1 2} & M_{1 3} & M_{1 4} \\
M_{2 1} & M_{2 2} & M_{2 3} & M_{2 4} \\
M_{3 1} & M_{3 2} & M_{3 3} & M_{3 4} \\
M_{4 1} & M_{4 2} & M_{4 3} & M_{4 4}
\end{array}
\right],
$$
where the matrices in the first column above are the following projections of the action of the differential:
\begin{itemize}
\item[$\bullet$] $M_{1 1}$ is the $(N_1+1) \times (N_1+1)$ matrix of columns $ \left[ \Re P_{N_1} \Pi_\phi  D \cRG(\phi_a,\psi_a) u_k \right]$,
\item[$\bullet$] $M_{3 1}$ is the $(N_1+1) \times (N_1+1)$ matrix of columns $ \left[ \Im P_{N_1} \Pi_\phi  D \cRG(\phi_a,\psi_a) u_k \right]$,
\item[$\bullet$] $M_{2 1}$ is the $(N_1+1) \times (N_2+1)$ matrix of columns
$ \left[ \Re P_{N_2} \Pi_\psi  D \cRG(\phi_a,\psi_a) u_k \right]$,
\item[$\bullet$] $M_{4 1}$ is the $(N_1+1) \times (N_2+1)$ matrix of columns
$ \left[ \Im P_{N_2} \Pi_\psi  D \cRG(\phi_a,\psi_a) u_k \right]$,
\end{itemize}
the matrices in the third column above are  the following projections of the action of the differential:
\begin{itemize}
\item[$\bullet$] $M_{1 3}$ is the $(N_1+1) \times (N_1+1)$ matrix of columns $ \left[ \Re P_{N_1} \Pi_\phi  D \cRG(\phi_a,\psi_a) \upsilon_k \right]$,
\item[$\bullet$] $M_{3 3}$ is the $(N_1+1) \times (N_1+1)$ matrix of columns $ \left[ \Im P_{N_1} \Pi_\phi  D \cRG(\phi_a,\psi_a) \upsilon_k \right]$,
\item[$\bullet$] $M_{2 3}$ is the $(N_1+1) \times (N_2+1)$ matrix of columns
$ \left[ \Re P_{N_2} \Pi_\psi  D \cRG(\phi_a,\psi_a) \upsilon_k \right]$,
\item[$\bullet$] $M_{4 3}$ is the $(N_1+1) \times (N_2+1)$ matrix of columns
$ \left[ \Im P_{N_2} \Pi_\psi  D \cRG(\phi_a,\psi_a) \upsilon_k \right]$,
\end{itemize}
the matrices in the second column above are  the following projections of the action of the differential:
\begin{itemize}
\item[$\bullet$] $M_{1 2}$ is the $(N_2+1) \times (N_1+1)$ matrix of columns $ \left[ \Re P_{N_1} \Pi_\phi  D \cRG(\phi_a,\psi_a) v_k \right]$,
\item[$\bullet$] $M_{3 2}$ is the $(N_2+1) \times (N_1+1)$ matrix of columns $ \left[ \Im P_{N_1} \Pi_\phi  D \cRG(\phi_a,\psi_a) v_k \right]$,
\item[$\bullet$] $M_{2 2}$ is the $(N_2+1) \times (N_2+1)$ matrix of columns
$ \left[ \Re P_{N_2} \Pi_\psi  D \cRG(\phi_a,\psi_a) v_k \right]$,
\item[$\bullet$] $M_{4 2}$ is the $(N_2+1) \times (N_2+1)$ matrix of columns
$ \left[ \Im P_{N_2} \Pi_\psi  D \cRG(\phi_a,\psi_a) v_k \right]$,
\end{itemize}
the matrices in the forth column above are  the following projections of the action of the differential:
\begin{itemize}
\item[$\bullet$] $M_{1 4}$ is the $(N_2+1) \times (N_1+1)$ matrix of columns $ \left[ \Re P_{N_1} \Pi_\phi  D \cRG(\phi_a,\psi_a) w_k \right]$,
\item[$\bullet$] $M_{3 4}$ is the $(N_2+1) \times (N_1+1)$ matrix of columns $ \left[ \Im P_{N_1} \Pi_\phi  D \cRG(\phi_a,\psi_a) w_k \right]$,
\item[$\bullet$] $M_{2 4}$ is the $(N_2+1) \times (N_2+1)$ matrix of columns
$ \left[ \Re P_{N_2} \Pi_\psi  D \cRG(\phi_a,\psi_a) w_k \right]$,
\item[$\bullet$] $M_{4 4}$ is the $(N_2+1) \times (N_2+1)$ matrix of columns
$ \left[ \Im P_{N_2} \Pi_\psi  D \cRG(\phi_a,\psi_a) w_k \right]$.
\end{itemize}


The Newton map $N_{\cRG,L_a}$ is now iterated until a good approximation $(\phi_0,\psi_0)$ of the renormalization fixed point is found. After that, a $2(N_1+1)+2(N_2+1)$ by $2(N_1+1)+2(N_2+1)$ matrix $L_0$ approximating the action of the differential $D \cRG(\phi_0,\psi_0)$ on $T_{(\phi_0,\psi_0)} \tilde{\cA}_1(U,V)$, is constructed.


This linear map $L_0$ is used in the following Theorem:

\begin{thm}\label{mainthmA} $\phantom{aaa}$\\
There exists a pair of polynomials $(\phi_0,\psi_0) \in \tilde{\cA}_1(U,V)$, where  $U=\D_{r_\phi}(c_\phi)$ and   $V=\D_{r_\psi}(c_\psi)$  are as in the Main Theorem 1, such that the operator $\cRG$ is analytic and compact in $B_r((\phi_0,\psi_0))$ in $\tilde{\cA}_1(U,V)$, $r=6.6610992 \cdot 10^{-11}$.

Furthermore, the Newton map $N_{\cRG,L_0}$ admits the following bounds:
\begin{itemize}
\item[$i)$] $\|(\phi_0,\psi_0)-N_{\cRG,L_0}(\phi_0,\psi_0)  \|_1 < \eps=1.33221983054516078 \cdot 10^{-11}$;
\item[$ii)$] $\| D N_{\cRG,L_0}(\phi,\psi)\| < \cD=0.287850876998644962$ for all $(\phi,\psi) \in B_\delta(  (\phi_0,\psi_0)  )$, $\delta=2.90241954016095172 10^{-11}$.
\end{itemize}
Therefore,
$$\epsilon < (1-\cD) \delta < -7.34735699562457612 \cdot 10^{-12}<0,$$
and the operator $N_{\cRG,L_0}$ has a fixed point in $B_\delta( (\phi_0,\psi_0) ) \subset \tilde{\cA}_1(U,V)$, while $\cRG$ has a fixed point  $(\phi_*,\psi_*)$ in  $B_r((\phi_0,\psi_0))$ in $\cA_1(U,V)$.
\end{thm}

\noindent
The compactness of $\cRG$ and $D|_{\zeta_*}\cRG$ follow from:
%The invariance of the fixed pair under renormalization allows one continue the domain of definition, analyticity and univalence of the maps $\phi_*$ and $\psi_*$.

\begin{prop}  \label{prop:renorm_compactness}
There exists $r'<r$ such that every  $(\phi, \psi) \in  B_{r'}((\phi_*,\psi_*))$  is in $\tilde{\cA}_1(\tU,\tV)$, where $\tU \supset \fD_{\tr_\phi} (\tc_\phi)$ and  $\tV \supset \fD_{\tr_\psi}(\tc_\psi)$, with $\tr_\phi=0.937$, $\tr_\psi=0.874$,  $\tc_\phi=0.6+i 0.09$ and $\tc_\psi=c_\psi$. 
\end{prop}







% \subsection{Limiting dynamics for the quadratic polynomial}


% In a neighborhood of $0$, the map 
% $$z \mapsto \bar{z}^\theta.$$
% conjugates $R_\theta^{q_n}$ to $R_\theta^{q_{n+1}}$ for all $n \ge N$. C. McMullen introduced in \cite{McM} the mapping
% $$\psi(z)=\phi^{-1}(\overline{\phi(z)}^\alpha),$$
% where $\phi$ is the conformal isomorphism of the unit disk with $\sd$, normalized such that $\phi(0)=0$, $0$ being the critical point of $\qp$, and proves  that there exists a neighborhood $U$ of $0$, and a constant $\epsilon$ such that  $\psi$ is well defined on $U \cap \overline{\sd}$, conjugates $\qp^{q_n}$ to $\qp^{q_{n+1}}$, and is $C^{1+\epsilon}$-conformal or anticonformal:
% $$\psi(z)= \lambda \overline{z} + O(|z|^{1+\epsilon}),$$
% here, $\lambda=\psi'(0)$ is the scaling ratio.

% The linearization of $\psi$ at $0$ will be called $\Lambda$:
% $$\Lambda(z)=\lambda \overline{z}.$$

% Among other results in \cite{McM}, McMullen proves the following 
% \begin{proposition}
% There exist two open topological disks $W_0 \subset \fC$ and $W_1 \subset \fC$, and two $\sigma$-proper branched coverings of $g_i: W_i \mapsto \fC$, $i=0,1$, such that
% $$g_i= \lim_{n \rightarrow \infty} \Lambda^{-n} \circ \qp^{q_n +i} \circ \Lambda^n,$$
% where the limit converges uniformly and exponentially on compact subsets of $W_i$.
% \end{proposition}


% To demonstrate that $(g_0,g_1)=\zeta^*$ it is, therefore, enough to show that there exists an iterate $n$, such that $\cRG^n (\qp^2,\qp) \in B_r((\phi_0,\psi_0))$. This is done by a straightforward computation.

% \begin{lemma}\label{lem:qpconvergence}
% $\cRG^{12}((\qp^2,\qp) \in B_r((\phi_0,\psi_0))$, and, therefore, $\cRG^n ((\qp^2,\qp) \rightarrow \zeta^*$. 
% \end{lemma}



\subsection{Existence of an unstable subspace}
Let $\lambda_*$ be as in the Main Theorem 1. 
We introduce an analytic operator
\begin{equation}
\cR_*(\eta,\xi)=\left(c \circ \lambda_*^{-1} \circ \eta \circ \xi \circ \lambda_* \circ c, c \circ \lambda_*^{-1} \circ \eta  \circ \lambda_*   \circ c \right),
\end{equation}
whose difference from $\cR$ is that the scaling factor $\lambda_*$ is fixed, rather than dependent on the pair $(\eta,\xi)$.
Of course, $\zeta_*$ is a fixed point of $\cR_*$. 

The following is easy to show:
\begin{lemma}
\label{difference-spectrum}
The spectra of the linear operators $D\cR_*^2|_{\zeta_*}$ and $D\cR^2|_{\zeta_*}$ satisfy
$$\text{sp}(D\cR_*^2|_{\zeta_*})=\text{sp}(D\cR^2|_{\zeta_*})\cup\{ 1\}.$$
\end{lemma}

Trivially, the map $\lambda_* \circ c$ where $\lambda_*=\lambda(\zeta^*)$ admits two invariant lines $l_+$ and $l_-$, which are perpendicular to each other.
 Specifically,
from  the invariance condition  $\lambda_* \circ c (x+a i x)=t (x+a i x)$ we have
\begin{eqnarray}
\nonumber \lambda_* \circ c (x+a i x)&=&\left( \Re(\lambda_*) x+ a \Im(\lambda_*) x \right)+i \left( \Im(\lambda_*) x- a \Re(\lambda_*) x \right)=t (x +a i x) \\
\nonumber &\implies& a \Re(\lambda_*) + a^2 \Im(\lambda_*)= \Im(\lambda_*) - a \Re(\lambda_*) \\
\nonumber &\implies&  a^2 \Im(\lambda_*) +2 a \Re(\lambda_*) - \Im(\lambda_*)=0 \\
\label{eq:apm} &\implies& a_\pm={ -\Re(\lambda_*) \pm |\lambda_*| \over \Im(\lambda_*)}.  
\end{eqnarray}
%Since $\Im(\lambda_*) \Re(\lambda_*) \ne 0$, $\lambda_* \circ c$ admits two distinct invariant lines, which will be denoted 
%Let us denote these two lines $l_+$ and $l_-$. $l_\pm(t)=t v_\pm$, $v_\pm$ being unit vectors in the direction of $l_\pm$. It is elementary to check that $\lambda_*(c(v_\pm))=\pm v_\pm$. The lines $l_+$ and $l_-$ are perpendicular to each other.


\begin{figure}[t]
\centering
% If not, use
\vspace{2mm}       % Give the correct figure height in cm
\includegraphics[height=7.0cm,width=6cm,angle=-90]{Inv_Lines.pdf}
\caption{Invariant lines: $l_+$ in blue and $l_-$ in red.}
\label{fig:lines} 
\end{figure}



Let $u(z)=\phi_u(z^2)$ and $v(z)=\psi_v(z^2)$ be analytic vector fields, such that $(\phi_u,\psi_v)$ is in $T_{(\phi_*,\psi_*)} \cA_1(U,V)$. 
Consider the action of the differential $D \cR_*(\zeta_*)(u,v)$
on a pair of such  fields   $(u_0,v_0)$, contained in some cone fields: $u_0(z) \in \cC_{u_0}(z)$ and $v_0(x) \in \cC_{v_0}(z)$:
\begin{equation}\label{eq:diraction}
(u_1,v_1):=(D \cR_*(\zeta_*)(u_0,v_0)= c \circ \lambda_*^{-1} \circ \left(u_0 \circ \xi_* + (\eta'_* \circ \xi_* ) \cdot v_0, u_0   \right) \circ \lambda_* \circ c.
\end{equation}

Clearly, $v_1 \in C_{v_0}(z)$, whenever $u_0(z) \in \cC_{u_0}(z)$,  if we set 
\begin{eqnarray} \label{eq:Cbetaw}
C_{v_0}(z):=c\left(\lambda_*^{-1}\left(C_{u_0}(\lambda_*(c(z)))\right)\right).
\end{eqnarray}
To find a pair of invariant cone fields, it is, therefore, sufficient to construct a cone field $C_{u_0}(z)$ invariant under the operator
\begin{equation}\label{eq:uoperator}
\cT(u):=c \circ \lambda_*^{-1} \circ \left(u \circ \xi_* \circ \lambda_* \circ c  +  (\eta'_* \circ \xi_* \circ \lambda_* \circ c ) \cdot (c \circ \lambda_*^{-1} \circ  u \circ |\lambda_*|^2) \right),
\end{equation}
and define $C_{v_0}(z)$ as in $(\ref{eq:Cbetaw})$. Notice that
\begin{eqnarray}
\label{eq:uoperator0} \cT(u)(0)&=&c \circ \lambda_*^{-1} \circ \left(u(\xi_*(0))  +  \eta'_*(\xi_*(0)) \cdot (c \circ \lambda_*^{-1} \circ  u(0) \right), \\
\label{eq:uoperator2n}\cT(u)(|\lambda|^{2n}) &=&c \circ \lambda_*^{-1} \circ(u(\xi_*(\lambda_*(c (|\lambda|^{2n}))))  +  \\
\nonumber &\phantom{=}& \phantom{c \circ \lambda_*^{-1}} + \eta'_*(\xi_*(\lambda_*(c(|\lambda|^{2n}) ))) \cdot (c \circ \lambda_*^{-1} \circ  u (|\lambda|^{2(n+1)}))),\\
\label{eq:uoperatorxi0} \cT(u)\circ \xi_*(0)&=&c \circ \lambda_*^{-1} \circ \left(u(|\lambda_*|^2) + \eta'_*(|\lambda_*|^2) \cdot (c \circ \lambda_*^{-1} \circ  u(|\lambda_*|^2) \right).
\end{eqnarray}

We will now proceed to construct a cone field $C_{u}$ which is invariant under the operator $\cT$ at points $|\lambda|^{2n}$, $n=0, \ldots, \infty$, and $0$. To that end, as formulas $(\ref{eq:uoperator0})-(\ref{eq:uoperatorxi0})$ suggest,  we will require bounds on derivatives of $\eta_*'$ at a sequence of points, as well as bounds on the images of the set $\lambda_*(c(|\lambda_*|^{2n}(Z)))$ under $\xi_*$.


\begin{lemma}\label{lem:conebounds}
The following bounds hold for the fixed point pair $\zeta_*$.
\begin{itemize}
\item[$1)$] $\eta_*'(\xi_*(0))=|\lambda_*|^{-2}$,
\item[$2)$] $|\arg(\eta_*'(|\lambda_*|^2)| <1.062$
\item[$3)$] $\sum_{n=1}^\infty |\arg(\eta_*'(\xi_*(\lambda_*(c(|\lambda_*|^{2n} )))))|<0.075=:\epsilon$,
\item[$4)$]  $| \arg(u(\xi_*(\lambda_*(|\lambda_*|^2))))-\arg(u(\xi_*(0)))   |<1.066$ and \\
$| \arg(u(\xi_*(\lambda_*(\bar{z}))))-\arg(u(\xi_*(0)))   |<0.398=\varepsilon$ for all $z \in [0,|\lambda_*|^4]$.
\end{itemize}
\end{lemma}
\begin{proof}

Notice that for the fixed point pair $\zeta_*=(\eta_*,\xi_*)$,
\begin{eqnarray}
\nonumber \eta_*' &=& c \circ \left( \eta'_* \circ \xi_* \cdot \xi_*'\right) \circ \lambda_* \circ c= c \circ \left( \eta'_* \circ \xi_* \cdot (c \circ \lambda_* \circ \eta_* \circ \lambda_* \circ c)'  \right) \circ \lambda_* \circ c\\
\label{eq:univbound} &=&\overline{\eta_*' \xi_* \circ \lambda_* \circ c} \cdot \left(\eta_*' \circ |\lambda_*|^2 \right)=\eta_*'  \implies  \overline{\eta_*' \xi_* \circ \lambda_* \circ c} = |\lambda_*|^{-2} { \phi_*'(z^2) \over \phi'_*(|\lambda_*|^4 z^2)}.
\end{eqnarray}
In particular, Claim $1)$ follows.
 
Consider the conformal transformation of $D_{\tr_\phi}(\tc_\phi)$, where $\tr_\phi$ and $\tc_\phi$ are as in Proposition $\ref{prop:renorm_compactness}$, onto the unit disk which maps the point $|\lambda_*|^4 z^2$ to zero:
$$ t^{-1}:=m \circ l, \quad l(x)={x  - \tc_\phi \over \tr_\phi}, \quad m(x):={x-l(|\lambda|^4 z^2) \over 1 - \overline{l(|\lambda|^4 z^2)} x}.$$
The function 
$$f(x)={ \phi_*(t(x)) -\phi_*(t(0))  \over \partial_x \phi_*(t(x))\arrowvert_{x=0} }$$
is univalent on a unit disk, normalized so that $f(0)=0$ and $f'(0)=1$, and
$$f'(t^{-1}(z^2))= { \phi_*'(z^2)   \over  \phi_*'(|\lambda|^4 z^2)  } {t'( (t^{-1}(z^2) ) \over t'(0)}.$$
We can now use  a classical distortion bound:
$$|\arg(f'(z))|<4 \arcsin{|z|},\ { \rm whenever} \ |z| <{1 \over \sqrt{2}},$$
in the class $\cS$ of functions $f$ univalent on $\fD$ and normalized so that $f(0)=0$, $f'(0)=1$, to get
\begin{eqnarray}
\nonumber && \left| \arg \left( { \phi_*'(z^2)   \over  \phi_*'(|\lambda|^4 z^2)  } {t'( (t^{-1}(z^2) )) \over t'(0) } \right) \right| = \\
\nonumber &&\phantom{aaaa}= \left| \arg \left( { \phi_*'(z^2)   \over  \phi_*'(|\lambda|^4 z^2)  } \right) +  \arg \left( \left( \tr_\phi^2-\overline{(|\lambda_*|^4 |z|^2 -\tc_\phi)}  (z^2-\tc_\phi) \right)^2 \right) \right| \\
\nonumber  &&\phantom{aaaa}= \left|\arg \left(t^{-1} (z^2) \right) \right| < 4 \arcsin\left( \left|\tr_\phi {1-|\lambda_*|^4 \over \tr^2_\phi-\overline{( |\lambda_*|^4 z^2-\tc_\phi  )}(z^2-\tc_\phi) }  |z|^2  \right|  \right).
\end{eqnarray}
Therefore,
\begin{eqnarray}
\nonumber | \arg{ \left( \eta_*'(\xi_*(\lambda_* c(z))) \right) }| &<&  4 \arcsin\left( \left|\tr_\phi{1-|\lambda_*|^4 \over \tr^2_\phi-\overline{( |\lambda_*|^4 z^2-\tc_\phi  )}(z^2-\tc_\phi) }  |z|^2  \right|  \right) \\
\label{eq:univbound1} &+& \left| \arg \left( \left( \tr_\phi^2-\overline{(|\lambda_*|^4 |z|^2 -\tc_\phi)}  (z^2-\tc_\phi) \right)^2 \right) \right|.
\end{eqnarray}
The bound in part $3)$ of the Lemma is obtain by evaluating first ten terms in the sum on a computer, and then estimating the remainder using simpler bounds on $\arcsin$ and $\arg$ functions in formula $(\ref{eq:univbound1})$. 
  
To demonstrate  the fourth claim, we first recall  the classical rotation bound
$$ \left|\arg{\left( { f(w) \over w      }   \right)}  \right|<\ln {1+|w| \over 1-|w|}$$
in the class $\cS$. In our case, $u=\varphi \circ q_2$  where  $\varphi$ is univalent on $\fD_{r_\phi}(c_\phi)$. Let $l(z)$ be the affine transformation that maps $\fD$ onto $\fD_{r_\phi}(c_\phi)$, set $a=l^{-1}(1)$, and let $M_\gamma(w)=(\gamma w +a )/(1+\bar{a} \gamma w)$, $|\gamma|=1$,   be the fractional linear transformation that preserves the unit disk and  maps $0$ into $a$. Consider the family 
$$f_\gamma(w)={ \varphi(l(M_\gamma(w)))-\varphi(1)  \over \varphi'(1) r_{\phi} \gamma (1-|a|^2)},$$
whose members are in $\cS$. Let point $w$ be the preimage in $\fD$  of some $\xi_*(\lambda_*(\bar{z}))$ under $l \circ M$, 
$$  \left| \arg{ \left( { f_\gamma(w) \over w }  \right)}  \right|=\left|\arg{\left( { \varphi(l(M_\gamma(w)))-\varphi(1)   \over  \varphi'(1) \gamma w} \right) } \right|<\ln {1+|w| \over 1-|w|}.$$ 
Now, for every $w \in \fD$ there is a choice of $\gamma=\gamma(w)$, such that  $ \varphi'(1) \gamma w \in \fR$, then
$$\left|\arg{ \left( \varphi(l(M_\gamma(w)))-\varphi(1) \right)  } \right|<\ln {1+|w| \over 1-|w|}.$$
A computer implementation of this bound results in the bounds in part $4)$.

Part $2)$ has been found directly on the computer.

\flushright $\square$
\end{proof}


We will now demonstrate existence of an invariant cone field $C$.



\begin{prop}\label{prop:unstable_dir} 
The spectral radii of the linear operators $D \cR_*^2(\zeta_*)$ and $D \cRG^2(\zeta_*)$ are larger than $1$. 
\end{prop} 
\begin{proof}
First, rewrite the equation $(\ref{eq:uoperatorxi0})$ as 
$$\cT(u)\circ \xi_*(0)=c \circ \lambda_*^{-1} \circ \left(u(|\lambda_*|^2) + c \circ \tilde{\lambda}_*^{-1} \circ  u(|\lambda_*|^2) \right), \quad \tilde{\lambda}_*= {\lambda_* \over \overline{\eta'_*(|\lambda_*|^2)}}.$$
Similarly to $(\ref{eq:apm})$, set  
$$\tilde{a}_\pm={ -\Re(\tilde{\lambda}_*) \pm |\tilde{\lambda}_*| \over \Im(\tilde{\lambda}_*)},$$
and let $\tilde{v}_\pm$ be the unit vectors of the invariant lines $\tilde{l}_\pm=(t+i \tilde{a}_\pm t)$ under the map $\tilde{\lambda}_* \circ c$.

Given a vector $v \in \fC$ and a positive $\alpha< \pi/2$, we denote a cone
$$C(v,\alpha)=\{x \in \fC \setminus \{0\}: |\arg(v)-\arg(x)| \le \alpha \} \cup \{0\}.$$

\medskip

\noindent  {\it Step 1)} We assume that cone field $\cC_u$ has some particular values at points $0$, $|\lambda_*|^2$ and $1$, and demonstrate that $\cT(\cC_u)(0) \subset {\rm Int} \ \cC_u(0)$ and $\cT(\cC_u)(1) \subset {\rm Int} \ \cC_u(1)$.


We make the following assumptions:
\begin{eqnarray}
\label{eq:conel2} \cC_u(|\lambda_*|^2)&:=&C(\tilde{v}_+,\alpha(|\lambda_*|^2)), \quad \alpha(|\lambda_*|^2)=\arccos(0.5967), \\
\label{eq:conel1} \cC_u(1)&:=&c(\lambda_*(C(\tilde{v}_+,\alpha(|\lambda_*|^2))))=:C(c(\lambda_*(\tilde{v}_+)),\alpha(1)))), \\
\label{eq:cone0} \cC_u(0)&:=&C(\tilde{v}_+,\alpha(0)), \quad \alpha(0)=\arccos(0.71).
\end{eqnarray}



The cone $C(\tilde{v}_+,\alpha(|\lambda_*|^2))$ is invariant under the map $c \circ \tilde{\lambda}_*^{-1}$, therefore for any $u(z)$ such that $u(|\lambda_*|^2) \in C(\tilde{v}_+,\alpha(|\lambda_*|^2))$, we have that  
$$\cT(u)\circ \xi_*(0) \in c(\lambda_*^{-1}(C(\tilde{v}_+,\alpha(|\lambda_*|^2))))=: C(c(\lambda_*^{-1}(\tilde{v}_+)),\alpha(1)),$$  
and, according to $(\ref{eq:uoperatorxi0})$, $\cT(u)(1) \in  \cC_u(1)$ whenever $u(|\lambda_*|^2) \in \cC_u(\lambda_*|^2)$. Additionally, the boundary $\partial \cC_u(|\lambda_*|^2)$ is a union of two lines, call them $l$ and $r$ and , and  the map $c \circ \tilde{\lambda}_*^{-1}$ interchanges them: $c \circ \tilde{\lambda}_*^{-1}(l)=r$ and   $c \circ \tilde{\lambda}_*^{-1}(r)=l$. Therefore, if $u(|\lambda_*|^2) \ne 0$ is in $\partial \cC_u(|\lambda_*|^2)=$, say $u(|\lambda_*|^2)\parallel l$, then $c \circ \tilde{\lambda}_*^{-1}(u(|\lambda_*|^2)) \parallel r$ and  $u(|\lambda_*|^2)+ c \circ \tilde{\lambda}_*^{-1}(u(|\lambda_*|^2)) \in {\rm Int} \ \cC_u(|\lambda_*|^2)$, and therefore, $\cT(\cC_u)(1) \in  {\rm Int} \ \cC_u(1)$.


We verify on the computer that  the cone $c(\lambda_*^{-1}(C(c(\lambda_*^{-1}(\tilde{v}_+)),\alpha(1))))$ is contained in  ${\rm Int} \ C(\tilde{v}_+,\alpha(0))$ (see Fig. \ref{fig:fields}). According to $(\ref{eq:uoperator0})$, this, together with the invariance of $\cC_u(0)$ under the map $c \circ \lambda_*^{-1} \circ c \circ \lambda_*^{-1}$ implies that $\cT(\cC_u)(0) \subset {\rm Int} \ \cC_u(0)$.


\begin{figure}
\centering
% If not, use
\vspace{2mm}       % Give the correct figure height in cm
\begin{tabular}{c c} 
\includegraphics[height=6cm,width=5.2cm,angle=-90]{Inv_Lines_1.pdf}  &  \includegraphics[height=6cm,width=5.2cm,angle=-90]{Inv_Lines_2.pdf}   \\
$a)$ & $b)$ \\
 \includegraphics[height=6cm,width=5.2cm,angle=-90]{Inv_Lines_3.pdf}  &  \includegraphics[height=6cm,width=5.2cm,angle=-90]{Inv_Lines_4.pdf}   \\
$c)$ & $d)$
\end{tabular}
\caption{The invariant lines $l_+$, $l_-$ and $\tilde{l}_+$ are drawn in blue, red and magenta, respectively.  $a)$ The cone $C(\tilde{v}_+,\alpha(|\lambda_*|^2))$ is between cyan and green lines.  b)  The cone  $C(c(\lambda_*(\tilde{v}_+)),\alpha(1)):=c(\lambda_*(C(\tilde{v}_+,\alpha(|\lambda_*|^2))))$ is between cyan and green lines. c) The image of the cone $C(c(\lambda_*(\tilde{v}_+)),\alpha(1))$  under $c \circ \lambda_*$ (orange to black) is contained in the cone $\cC_u(0)$ (green to cyan). d) $\varepsilon$-thickening (yellow and violet) of the  image of the cone $C(c(\lambda_*(\tilde{v}_+)),\alpha(1))$  under $c \circ \lambda_*$ (orange to black) is, again,  contained in the cone $\cC_u(0)$ (green to cyan)}
\label{fig:fields}       % Give a unique label
\end{figure}



\medskip

\noindent  {\it Step 2)} We continue to make assumptions on the field $\cC(z)$:
\begin{equation}
\label{eq:conesdef} \cC_u(|\lambda_*|^{2 n})=C(\tilde{v}_+,\alpha(0)+  \sum_{i=n}^\infty |\arg(\eta_*'(\xi_*(\lambda_*(c(|\lambda_*|^{2i} )))))|), \quad n=2,3, \ldots .
\end{equation}

By the bound in part 4) in Lemma \ref{lem:conebounds},  $u(\xi_*(\lambda_*|\lambda_*|^{2n}))$, $n=2,3, \ldots$ is in $C(c(\lambda_*(\tilde{v}_+)) , \alpha(1)+\varepsilon)$, $\varepsilon=0.398$ as in part 4) of Lemma \ref{lem:conebounds}, whenever $u(\xi_*(0)) \in C(c(\lambda_*(\tilde{v}_+)),\alpha(1))$. We verify on the computer that 
\begin{equation}\label{eq:cone1}
c(\lambda_*^{-1}(C(c(\lambda_*(\tilde{v}_+)),\alpha(1)+\varepsilon))) \subset {\rm Int} \ \cC_u(0):= {\rm Int} \ C(\tilde{v}_+,\alpha(0))
\end{equation}
(see Fig. \ref{fig:fields}). Furthermore,
\begin{eqnarray}
\nonumber  c \circ \lambda_*^{-1} \! \! & \! \! \circ  \!  \! & \! \! \left( \eta'_*(\xi_*(\lambda_*(c(|\lambda|^{2n}) )) \cdot (c \circ \lambda_*^{-1} \circ  u (|\lambda|^{2(n+1)})) \right) \in \\ 
\label{eq:cone2}  \! \! &\! \! \in \! \! & \! \! C(\tilde{v}_+,\alpha(0)+\sum_{i=n}^\infty |\arg(\eta_*'(\xi_*(\lambda_*(c(|\lambda_*|^{2i} )))))|)
\end{eqnarray}
whenever  
$$u (|\lambda|^{2(n+1)}) \in C(\tilde{v}_+,\alpha(0)+\sum_{i=n+1}^\infty |\arg(\eta_*'(\xi_*(\lambda_*(c(|\lambda_*|^{2i} )))))|).$$
Conditions $(\ref{eq:cone1})$ and $(\ref{eq:cone2})$ imply that $\cT(u)(|\lambda_*|^{2 n})$, given by $(\ref{eq:uoperator2n})$,  is in ${\rm Int} \ \cC_u(|\lambda_*|^{2 n})$, as defined in $(\ref{eq:conesdef})$ whenever $u(|\lambda_*|^{2 (n+1)}) \in \cC_u(|\lambda_*|^{2 (n+1)})$, $n=1,2, \ldots$. Since $\arccos(0.5967) - \arccos(0.71)>0.075$, the cone  $\cT(\cC_u)(|\lambda_*|^{2})$, according to $(\ref{eq:cone2})$, is contained in ${\rm Int} \ \cC_u(|\lambda_*|^2)$ as defined in $(\ref{eq:conel2})$. This demonstrates invariance of cones $\cC(|\lambda_*|^{2 n})$ for all $n=1,2, \ldots$.

\medskip


%%%%%%%%%%%%%%%%%%%%%%%%%%%%%%%%%%%%%  begin ignore
%%%%%%%%%%%%%%%%%%%%%%%%%%%%%%%%%%%%%
\ignore{
%%%%%%%%%%%%%%%%%%%%%%%%%%%%%%%%%%%%%
%%%%%%%%%%%%%%%%%%%%%%%%%%%%%%%%%%%%%
\noindent  {\it Step 3)}  We will now demonstrate that the operator $D \cR_*(\zeta_*)$ has an eigenvalue outside the unit circle which corresponds to an eigenvector which is not tangent to the manifold of almost commutative pairs, that is, not an eigenvector in $T_{\zeta_*} \cM(U,V)$.

The equation $(\ref{eq:uoperator0})$ together with part $1)$  of Lemma $\ref{lem:conebounds}$, and the fact that $\alpha(0)<\pi/4$,  implies that 
$$|\cT(u)(0)|  \ge  {1 \over |\lambda_*|^4} |u(0)|.$$
In particular, this implies that $\cT$ maps a vector field $u$ which are non-zero in a neighborhood of point $0$ into a field which is non-zero on some neighborhood of $0$, and which, furthrmore, there exists a neighborhood $\tilde{Z}$ of $0$ and $1 < \mu_u  \le {|\lambda_*|^{-4}}$, such that
%whenever $u(z) \in \cC_u(z)$ is non-zero in a neighborhood of $0$. In particular, this implies that $D \cR_*(\zeta_*)$ maps a pair $(u,v)$ of vector fields, which are non-zero in a neighborhood of point $0$ into a pair which is non-zero on some neighborhood of $0$, and which, furthrmore, is stretched in a neighborhood of $0$: there exists a neighborhood $\tilde{Z}$ of $0$ and $1 < \mu_u  \le {1 \over |\lambda_*|^4}$, such that
\begin{equation}
\label{eq:expansion}  |\cT(u)(z)| \ge \mu_u |u(z)|
\end{equation}
for all $z \in \tilde{Z}$.

By a standard result (see, e.g. Proposition 6.2.12 and Corollary 6.2.13 in \cite{KH}), expansion of non-zero vectors in $\cC_u(z)$ for $z \in \tilde{Z}$   and   the   invariance condition $\cT(\cC_u) \subset {\rm Int} \ \cC_u$ on the set $\tilde{Z} \cap E$, where $E=\{|\lambda_*|^{2n}\}_{n=0}^\infty \cup 0$, implies that any  analytic direction field in $\cC_u(z)$ converges under the operator $\cT$ in $\tilde{Z} \cap E$ to a  analytic invariant direction field for $\cT$.  Since the set $\tilde{Z} \cap E$ contains a limit point $0$,  any  analytic direction field in $\cC_u(z)$ in $Z$ converges under the operator $\cT$ in all of $Z$. 


 % Indeed, let $u$ and $w$ be any two analytic non-zero vector fields in $Z$, such that $u(|\lambda_*|^{2 n}), w(|\lambda_*|^{2 n}) \in \cC(|\lambda_*|^{2 n})$, $n=0,1, \ldots$ and $u(0), w(0) \in \cC_u(0)$,  Then, since the property of being non-zero in a neighborhood of $0$ is preserved by $\cT$,  both sequences $\arg\left(\cT^n(u)(z)/\|\cT^n(u)\|_1\right)$ and $\arg\left( \cT^n(w)(z)/\|\cT^n(w)\|_1 \right)$  are sequences of bounded real-analytic functions on some neighborhoods $\tilde{Z} \subset Z$ and $W \subset W$ of $0$, and converge uniformly on $Z'$ and $W'$. The functions $a_\infty(z)=\lim_{n \rightarrow \infty}\arg\left( \cT^n(u)(z)/\|\cT^n(u)\|_1 \right)$ and $b_\infty(z)=\lim_{n \rightarrow \infty}\arg\left( \cT^n(w)(z)/\|\cT^n(w)\|_1 \right)$, real-analytic on $Z'$ and $W'$, respectively,  coincide on a set with a limit point, and therefore, have to coincide everywhere on $Z' \cup W'$.


The equation $(\ref{eq:expansion})$ implies the claim about the spectrum. Since $D \cR_*^2(\zeta_*)$ is a compact operator whose spectrum, outside of $0$, consists of eigenvalues, the largest eigenvalue is larger than $\mu_u$. 

Notice, that this eigenvector $(u_*,v_*)$ can  not lie in $T_{\zeta_*}\cM(U,V)$ since both $u_*(0)$ and $v_*(0)=c \circ \lambda^{1}_* u_*(0)$ are non-zero.

It is easy to show that the spectrum of the operator $\cR_*$ differs from that of $\cR$ by a single neutral eigenvalue $\delta=1$ (see Subsection \ref{subsec:coord_changes}). In particular, the expanding parts of the spectra coincide.

%%%%%%%%%%%%%%%%%%%%%%%%%%%%%%% end ignore
%%%%%%%%%%%%%%%%%%%%%%%%%%%%%%%%%%%%%%%%%%
}
%%%%%%%%%%%%%%%%%%%%%%%%%%%%%%%%%%%%%%%%%%

\medskip

\noindent  {\it Step 3)}  Assume now, that $u(z)=a z^2+b z^4 + \ldots \in \cC(z)$ is non-zero in some punctured neighborhoods $U \setminus \{ 0\}$ of $0$ and  $V \setminus \{1\}$ of $1$. Then, the equation $(\ref{eq:uoperator0})$ together with part $1)$  of Lemma $\ref{lem:conebounds}$, and the fact that $\alpha(|\lambda_*|^{2 n})<\pi/4$,  implies that 
\begin{equation}
\label{eq:cT2n} |\cT(u)(|\lambda_*|^{2 n})| > {1 \over |\lambda_*|^4} |u( |\lambda_*|^2 |\lambda_*|^{2n})| + \gamma {1 \over |\lambda_*|}  |u(\xi_*(\lambda_* |\lambda_*|^{2 n}))|,
\end{equation}
for some $\gamma>0$. Since $u(z)=O(z^2)$ in a neighborhood of $0$ and $u(1+z)=O(z)$ in a neighborhood of $1$, there exists $C>0$ such that $\left||u(z)|-{1 \over 2} |u''(0)||z|^2 \right|<C |z|^4$ and $\left| |u(\xi_*(z))|-2|u'(1)||\psi'(0)||z| \right|<C |z|^2$ for sufficiently small $z$. Therefore, according to $(\ref{eq:cT2n})$, for sufficiently large $n$,
\begin{eqnarray}
\nonumber |\cT(u)(|\lambda_*|^{2 n})| \!\! & \!\!> \!\!&  \!\! \left\{ \! {1 \over 2} |u''(0)||\lambda_*|^{4 n} \! \!-C |\lambda_*|^4 |\lambda_*|^{8 n} \! \right\}\!+\! \gamma \left\{ 2|u'(1)||\psi'(0)||\lambda_*|^{2 n} \!\!-C |\lambda_*||\lambda_*|^{4 n}   \right\}, \\
\nonumber  \!\!& \!\!= \!\!&  \!\!\left\{ {1 \over 2} |u''(0)||\lambda_*|^{4 n} +C |\lambda_*|^{8 n}  \right\}+ \gamma \left\{ 2|u'(1)||\psi'(0)||\lambda_*|^{2 n} -C |\lambda_*||\lambda_*|^{4 n} - \right. \\
\nonumber  \!\!& \!\!\phantom{=} \!\!&  \!\!\phantom{\left\{ {1 \over 2} |u''(0)||\lambda_*|^{4 n} +C |\lambda_*|^{8 n}  \right\}}- \left. C(1+|\lambda_*|^4)  |\lambda_*|^{8 n} \right\} \\
\nonumber  \!\!& \!\!> \!\!& \!\! |u(|\lambda_*|^{2 n})|\!+\! \gamma \tilde{C} |\lambda_*|^{2 n},
\end{eqnarray}
for some constant $\tilde{C}>0$. Therefore, $\cT$ strictly increases the uniform norm of $u$ on small neighborhoods of $0$.  The corresponding statement in the norm $||\cdot ||_1$ follows from \propref{norm-equiv}.


We remark, that $D \cR_*(\zeta_*)(u,v)= D \cR(\zeta_*)(u,v)$ whenever $v(0)=0$. Furthermore for such $(u,v)$, we have on a neighborhood of $0$
$$D \cRG(\zeta_*)(u,v)(z)=(\tilde{u},\tilde{v})(z)+(D a(\tilde{u},\tilde{v}) z^4 + D b(\tilde{u},\tilde{v}) z^6,0),$$
where $(\tilde{u},\tilde{v})=D \cR_*(\zeta_*)(u,v)$. Therefore, for the operator $\cT$ is modified as follows:
$$\cT(u)(z):=c \circ \lambda_*^{-1} \circ \left(u(\xi_*(\lambda_*(c(z))))  +  \eta'_*(\xi_*(\lambda_*(c(z)))) \cdot (c \circ \lambda_*^{-1} \circ  u(|\lambda_*|^2 z )) \right) +O(z^4),
$$
and the constructed cones $\cC_u(|\lambda_*|^{2 {n}}) \times \cC_v(|\lambda_*|^{2 n})$ and $\cC_u(0) \times \cC_v(0)$ remain invariant  under the action of the differential $D \cRG(\zeta_*)$ for large $n$ (small neighborhood of zero). An analytic direction field in a punctured neighborhood of zero converges on this neighborhood under the action of $\cT$, and, hence, converges everywhere on $Z$.
\flushright $\square$
\end{proof}




% Furthermore, consider a family $\zeta^t_*=(\eta^t_*,\xi^t_*)=(\eta_*+t u_*, \xi_* + t v_*)$, where $(u_*,v_*)$ is an expanding eigenvector of  $\cR$.  Notice that $\zeta^t_*$ is tangent to the almost commuting submanifold $\cM(Z,W)$:
%$$(\eta^t_* \circ \xi^t_*)^{(n)}(0)-(\xi^t_* \circ \eta^t_*)^{(n)}(0)=\left( t (u_* \circ \xi_* + ( \eta'_* \circ \xi_*(0) ) \cdot v_*$$



%  $n=0,1,2$, then so is the pair $\tilde{\zeta}_t:=\cR(\zeta_t)$. Indeed, according to $(\ref{eq:rcommute0})$ 
% \begin{eqnarray}
% \nonumber (\tilde{\eta}_t \circ \tilde{\xi}_t)(0)-(\tilde{\xi}_t \circ \tilde{\eta}_t)(0)&=& \left(c \circ \lambda^{-1} \circ \eta_t \circ \xi_t  \circ \eta_t \right)(0)-\left(c \circ \lambda^{-1} \circ \eta_t  \circ \eta_t \circ \xi_t \right)(0) \\
% \label{eq:tangency0}&=&  \left(c \circ \lambda^{-1} \circ \eta_t \circ \xi_t  \circ \eta_t \right)(0)-\left(c \circ \lambda^{-1} \circ \eta_t  \circ \xi_t \circ \eta_t \right)(0)+O(t^2)=O(t^2),
% \end{eqnarray}
% and similarly for equations $(\ref{eq:rcommute1})$ and $(\ref{eq:rcommute2})$. Therefore, $D \cR(\zeta_*)$ maps ${T_{\zeta_*} \cM(Z,W)}$ to itself, and $\sigma(D \cR(\zeta_*) \arrowvert_{T_{\zeta_*} \cM(Z,W)})= \sigma(D \cRG(\zeta_*) \arrowvert_{T_{\zeta_*} \cE(Z,W)})$. Therefore, to demonstrate that the spectral radius of $D \cRG(\zeta_*)$ is larger than $1$, it is sufficient to show that there exist a vector field $(u,v)$ in ${T_{\zeta_*} \cM(Z,W)}$ with values in $C(z) \times c(\lambda_*^{-1}(C(\lambda_* \bar{z})))$.

% \begin{proof}
% First, we write out the condition  $(\eta_t \circ \xi_t)(0)-(\xi_t \circ \eta_t)(0)=O(t^2)$ of tangency  in more detail:
% \begin{eqnarray}
% \nonumber O(t^2)&=&(\eta_t \circ \xi_t)(0)-(\xi_t \circ \eta_t)(0)  \\
% \nonumber &=&   \eta_* \circ \xi_* (0) + t u \circ \xi_*(0) + t \eta'_* \circ \xi_*(0) \cdot v(0)-\xi_* \circ \eta_* (0) - t v \circ \eta_*(0) - t \xi'_* \circ \eta_*(0) \cdot u(0) + O(t^2) \implies \\
% \label{eq:tangency00}  &\implies&  u \circ \xi_*(0) + \eta'_* \circ \xi_*(0) \cdot v(0) - v \circ \eta_*(0)-  \xi'_* \circ \eta_*(0) \cdot u(0) = 0.
% \end{eqnarray}

% Consider a vector field $(u,v) \in T_{\zeta_*} \cM(Z,W)$,  such that $u(|\lambda_*|^{2 n})=a_n \tilde{v}_+$, $n=1,2, \ldots$,   $u(1)=a_0 c \circ \lambda_*^{-1} \circ  \tilde{v}_+$, and  $v(\bar{\lambda}_*^{-1} |\lambda|^{2n})=b_n c \circ \lambda_*^{-1} \tilde{v}_+$,  $n=0,1, \ldots$,     where $a_n$ and $b_n$  are some  real positive numbers. The condition $(\ref{eq:tangency00})$ becomes for these vector fields
% \begin{eqnarray}
%  a_0  c \circ \lambda_*^{-1} \circ \tilde{v}_+  + |\lambda_*|^{-2}  b_\infty c \circ \lambda_*^{-1} \tilde{v}_+ - b_1  c \circ \lambda_*^{-1} \circ  \tilde{v}_+ -  \xi'_* \circ \eta_*(0) \cdot a_\infty \tilde{v}_+ = 0.
% \end{eqnarray}
% Notice, that the commutation condition $k$ implies
% $$\xi'_* \circ \eta_*(0) \phi_*(0)  \phi'_*(0)= \eta_*' \circ \xi_*(0) \psi_*(0) \psi'_*(0)= |\lambda_*|^{-2}   \psi'_*(0)  \implies \xi'_* \circ \eta_*(0)=|\lambda_*|^{-2} \lambda_*^{-1}  {\psi'_*(0) \over \phi'_*(0)}= \lambda_*^{-2} {\overline{\phi'_*(0)} \over \phi'_*(0)}$$
% \end{proof}


  







\subsection{Renormalization hyperbolicity}


To diagonalize the differential, we first find numerically eigenvalues and eigenvectors of the matrix $L_0$. It is elementary to check that the anti-linearity of the renormalization differential implies that the eigenvalues of $L_0$ come in pairs: if $\lambda$ is an eigenvalue, then so is $-\lambda$:
\begin{lemma} \label{lem:spectra}
Let $\sigma$  be the spectrum of the linear operator $D \cRG[(\phi_*,\psi_*)] \circ c$ on $T \cA_1(U,V)$. Then the spectrum of the linear operator $D \cRG[(\phi_*,\psi_*)]$ in $T \tilde{\cA}_1(U,V)$ is 
$$\tilde{\sigma}=-\sigma \cup \sigma.$$
\end{lemma} 
We verify numerically, that the matrix $L_0$ has $2(N_1+1)+2(N_2+1)$ eigenvalues, some of which have a multiplicity larger than $1$. We compute all real  eigenvalues, together with their eigenvectors, and the complex pairs, together with their eigenplanes, numerically using the standard linear algebra procedure in ADA. 


At the next step, we choose the first $E \equiv 2(E_1+1)+2(E_2+1)$ eigenvectors $e_i$, $E=30$ and $E_2=20$, as the basis in the following linear subspace:
$${\rm span} \left\{ \{ (l_\phi^k,0) \}_{k=0}^{E_1},\{(0,l_\psi^k)\}_{k=0}^{E_2},\{(i l_\phi^k,0)\}_{k=0}^{E_1},\{(0,i l_\psi^k)\}_{k=0}^{E_2} \right\},$$ 
keeping 
$$\left\{ \{ (l_\phi^k,0) \}_{k=E_1+1}^{N_1},\{(0,l_\psi^k)\}_{k=E_2+1}^{N_2},\{(i l_\phi^k,0)\}_{k=E_1+1}^{N_1},\{(0,i l_\psi^k)\}_{k=E_1+1}^{N_2} \right\}$$ 
as the remaining basis vectors. The new basis will be referred to as $\cE$. 

To find the linear transformation from the basis $\cL$ to $\cE$, we first find the inverse $B^{-1}$ the matrix $B=\left[ \cE \right]$ numerically, and then find a bound on the inverse using a Newton method in the equation $B B^{-1}=\I$. 

Finally, we consider a projection of the bound $D$ on operator $D \cRG$ in  $B_r$ on the complimentary subspace to the span of the approximate eigenvectors $e_0$ and $e_1$ which numerically span the expanding subspace:
$$T=(I-S_0 -S_1) D (\I -S_0 -S1), \quad S_i=\left[
\begin{array}{c}
e_i^0  B^{-1}_i \cr
e_i^1  B^{-1}_i \cr
\ldots   \cr
e_i^N  B^{-1}_i \cr
\end{array}
\right], 
$$
where $e_i^k$ are the $k$-th component of the $i$-th eigenvector, and $B^{-1}_i$ stands for the $i$-th row of the matrix $B$. We compute rigorously that
$$\| T^3 \| <1,$$ 
which implies that the stable subspace of $D \cRG[(\phi_*,\psi_*)]$ in $T_{(\phi^*,\psi^*)} \tilde{A}_1(U,V)$ has real codimension two.  This together with Proposition \ref{prop:unstable_dir} implies Part (v) of the Main Theorem 1.  


%%%%%%%%%%%%%%%%%%%%%%%%%%%%%% begin ignore
%%%%%%%%%%%%%%%%%%%%%%%%%%%%%%
\ignore{
%%%%%%%%%%%%%%%%%%%%%%%%%%%%%%
\subsection{Renormalization spectrum associated with the coordinate changes}\label{subsec:coord_changes}
 
In this section we will consider the part of the renormalization spectrum associated with the coordinate changes. Let $\sigma$ be a generator of a coordinate transformation, then
\begin{equation} \label{eq:change} (\id +t \sigma)^{-1} \circ \eta \circ (\id +t \eps \sigma)=\eta+t h_{\eta,\sigma}+O(t^2), \ {\rm where} \ h_{\eta,\sigma} = -\sigma \circ \eta + \eta' \cdot \sigma.
\end{equation}

Consider the operator  
$$\cR_* (\eta,\xi)=(\las^{-1} \circ \eta \circ \xi \circ \las, \las^{-1} \circ \eta \circ \las),$$
where $\las=-|\eta_*(0)|$. Then
\begin{eqnarray}
\nonumber D \cR_*(\eta_*,\xi_*) \left(h_{\eta_*,\sigma},h_{\xi_*,\sigma}\right)&=&( \las^{-1} \circ \left(-\sigma \circ \eta_* \circ \xi_*  +(\eta_* \circ \xi_*)' \cdot \sigma  \right) \circ \las, \\
\nonumber &\phantom{=}&\phantom{(} \las^{-1} \circ \left(-\sigma \circ \eta_*  +\eta_*' \cdot \sigma  \right) \circ \las ) \\
\nonumber &=&( - \las^{-1} \circ \sigma \circ \las \circ \eta_*  +\eta_*' \cdot (\las^{-1} \sigma  \circ \las), \\
\nonumber &\phantom{=}&\phantom{(} -\las^{-1} \circ \sigma \circ \las \circ \xi_*  +\xi_*' \cdot (\las^{-1} \circ \sigma  \circ \las) ) \\
\nonumber &=&  \left(h_{\eta_*,\las^{-1} \circ \sigma \circ \las},h_{\xi_*, \las^{-1} \circ \sigma \circ \las} \right),
\end{eqnarray}
and is equal to $\kappa \left(h_{\eta_*,\sigma},h_{\xi_*,\sigma}\right)$ whenever 
$$\las^{-1} \circ \sigma \circ \las= \kappa \sigma.$$
Consider monomial generators $\sigma_k(z)=z^k$, corresponding to the eigenvectors
$$e_k=\left(h_{\eta_*,\sigma_k},h_{\xi_*,\sigma_k}\right),$$ 
then
$$\las^{-1} \circ \sigma_k \circ \las= \las^{k-1} \sigma_k, $$
In particular, the constant generator of translations, $\sigma_0(z)=1$, corresponds to an expanding eigenvector of the operator $D \cR_*$, while the generator of rescalings $\sigma_1(z)=z$ corresponds to a neutral eigendirection of $D \cR_*$. At the same time $\span\{e_1\}$ is the null eigenspace for the operator $D \cR$ in $T_{(\eta_*,\xi_*)}\cE(U,V)$.

Perturbations in the  directions of the eigenvectors $e_k$, $k \ge 1$, preserve the class of a.c.s. pairs. Indeed, let $\tilde{\eta}=(\id +\eps \sigma_k)^{-1} \circ \eta \circ  (\id +\eps \sigma_k)$ and $\tilde{\xi}=(\id +\eps \sigma_k)^{-1} \circ \xi \circ  (\id +\eps \sigma_k)$, then $(\ref{eq:accond})$ implies
\begin{eqnarray}
\nonumber \tilde{\eta} \circ \tilde{\xi}(0)&=& -\sigma_k \circ \eta \circ \xi(0)  + (\eta \circ \xi) '(0) \cdot \sigma_k(0) +O(\eps^2)\\
\nonumber &=&-\sigma_k \circ \xi \circ \eta(0)  + (\xi \circ \eta) '(0) \cdot \sigma_k(0) +O(\eps^2)\\
\nonumber &=&\tilde{\xi} \circ \tilde{\eta}(0)+O(\eps^2) \\
\nonumber (\tilde{\eta} \circ \tilde{\xi})'(0)&=& -\sigma_k' \circ \eta \circ \xi(0) \cdot  (\eta \circ \xi)'(0)    + (\eta \circ \xi) ''(0) \cdot \sigma_k(0)+ (\eta \circ \xi) '(0) \cdot \sigma_k'(0)+\\
\nonumber &\phantom{=}& \phantom{-\sigma_k' \circ \eta \circ \xi(0) \cdot  (\eta \circ \xi)'(0)    + (\eta \circ \xi) ''(0) \cdot \sigma_k(0)} + O(\eps^2) \\
\nonumber &=&(\tilde{\xi} \circ \tilde{\eta})'(0)+O(\eps^2),
\end{eqnarray}
\begin{eqnarray}
\nonumber (\tilde{\eta} \circ \tilde{\xi})''(0)&=& -\sigma_k'' \circ \eta \circ \xi(0) \cdot  ((\eta \circ \xi)'(0))^2-\sigma_k' \circ \eta \circ \xi(0) \cdot  (\eta \circ \xi)''(0)   \\
\nonumber &\phantom{=}&+ (\eta \circ \xi) '''(0) \cdot \sigma_k(0)+2 (\eta \circ \xi) ''(0) \cdot \sigma_k'(0)+ (\eta \circ \xi) '(0) \cdot \sigma_k''(0)+O(\eps^2)\\
\nonumber &=&(\tilde{\xi} \circ \tilde{\eta})''(0)+O(\eps^2), \\
\nonumber (\tilde{\eta} \circ \tilde{\xi})'''(0)&=& -\sigma_k''' \circ \eta \circ \xi(0) \cdot  ((\eta \circ \xi)'(0))^3-3 \sigma_k'' \circ \eta \circ \xi(0) \cdot (\eta \circ \xi)'(0) \cdot  (\eta \circ \xi)''0) -\\\nonumber &\phantom{=}& -\sigma_k' \circ \eta \circ \xi(0) \cdot  (\eta \circ \xi)'''(0)  + (\eta \circ \xi)^{(4)}(0) \cdot \sigma_k(0) +\\
\nonumber &\phantom{=}& +3 (\eta \circ \xi) '''(0) \cdot \sigma_k'(0)+3 (\eta \circ \xi) ''(0) \cdot \sigma_k''(0)+\\
\nonumber &\phantom{=}& +(\eta \circ \xi) '(0) \cdot \sigma_k'''(0) +O(\eps^2) \\
\nonumber &=&(\tilde{\xi} \circ \tilde{\eta})'''(0)+\left\{ \left( (\eta \circ \xi)^{(4)}(0)-(\xi \circ \eta)^{(4)}(0) \right) \cdot \sigma_k(0) \right\}+O(\eps^2). 
\end{eqnarray}
In the last expression in parenthesis $\sigma_k(0)=0$ for all $k > 0$, but is non-zero for $k=0$. Therefore, the eigenvector $e_0$, corresponding to translations of coordinates, is the only eigendirection of $D \cR_*$ in $T_{(\eta_*,\xi_*)}\cE(U,V)$  associated with coordinate changes which does not lie in the tangent space $T_{(\eta_*,\xi_*)}\cM(U,V)$ at $(\eta_*,\xi_*)$ to the manifold of a.c.s. pairs.

%%%%%%%%%%%%%%%%%%%%%%%%%% end ignore
}
%%%%%%%%%%%%%%%%%%%%%%%%%%%%%%%%%%%%%
